\documentclass[a4paper, 12pt]{article}
\usepackage[english, russian]{babel}
\usepackage[T2A]{fontenc}
\usepackage[utf8]{inputenc}
\usepackage{graphicx}
\usepackage{amsmath}
\usepackage{amssymb}
\usepackage{amsfonts}
\usepackage{array}
\usepackage{booktabs}
\usepackage{wrapfig}
\usepackage{caption} %заголовки плавающих объектов
\captionsetup[table]{justification=centering} %установки для заголовков таблиц
\captionsetup[figure]{justification=centering} %установки для заголовков таблиц


\begin{document}
	\begin{center}
		\subsubsection*{Моделирование стационарного распределения температур в мощных полупроводниковых лазерах AlGaAs/GaAs/InGaAs}
		Куликовских Илья, СПбПУ группа 5030302/30801
		
		Руководитель: Корнышов Григорий, ФТИ им. Иоффе, лаборатория физики полупроводниковых гетероструктур.
	\end{center}
	
	\quad В данной работе исследуется распределение температуры в торцевом полупроводниковом лазере на основе гетероструктуры AlGaAs/GaAs с квантовыми ямами InGaAs. Работа проведена в лаборатории физики полупроводниковых гетероструктур Физико-технического института им. А. Ф. Иоффе.
	
	\quad Ограничение мощности современных полупроводниковых лазеров связано с разогревом полупроводниковой гетероструктуры. Тепло выделяется как при протекании тока из-за омического сопротивления через слои гетероструктуры, так и из-за безызлучательной рекомбинации в активной области. Разогрев приводит к росту внутренних потерь и снижению дифференциальной эффективности, а в конечном итоге к катастрофическому разрушению прибора. Построение теоретической модели распределения тепла от тока накачки в мощном полупроводниковом лазере, работающем в непрерывном режиме работы, позволит лучше прогнозировать его предельные характеристики. Полученные результаты будут использованы для оптимизации толщин слоёв гетероструктуры новых мощных лазеров. 
	
\end{document}